\documentclass[11pt,a4paper]{ctexart}

\usepackage{amsmath,amssymb,mathtools}
\usepackage{geometry}
\usepackage{booktabs}
\usepackage{hyperref}
\usepackage{xcolor}
\usepackage{enumitem}

\geometry{margin=2.4cm}
\hypersetup{
  colorlinks=true,
  linkcolor=blue!60!black,
  urlcolor=blue!60!black,
  citecolor=blue!60!black
}
\setlist[itemize]{leftmargin=2em}
\setlist[enumerate]{leftmargin=2em}

\title{World Action Models 的三种典型范式}
\author{}
\date{\today}

\begin{document}
\maketitle

\section{总览}

World Action Model，简称 WAM，试图把世界预测与动作生成统一起来。它不同于普通策略模型只学习
\[
  p(a_t\mid o_t,\ell),
\]
也不同于纯 world model 只学习
\[
  p(o_{t+1:t+H}\mid o_{\le t},a_{t:t+H-1}).
\]
WAM 的核心是建模动作与未来世界之间的耦合关系。典型目标可以写为：
\[
  p_\theta(o_{t+1:t+H},a_{t:t+H-1}\mid o_{\le t},\ell).
\]

当前 WAM 文献中可以抽象出三种代表性范式：
\begin{enumerate}
  \item Joint Video-Action Modeling：联合生成未来视频和动作；
  \item Imagine-Then-Execute：先想象未来，再根据未来执行；
  \item Fast-WAM：训练时学习未来预测，推理时跳过显式未来生成。
\end{enumerate}

\section{范式一：Joint Video-Action Modeling}

\subsection{核心思想}

Joint Video-Action Modeling 把未来观测和动作放在同一个生成建模问题中。模型不是先生成视频再生成动作，也不是只预测动作，而是直接学习联合分布：
\[
  p_\theta(o_{t+1:t+H},a_{t:t+H-1}\mid o_{\le t},\ell).
\]

直观上，模型需要同时回答两个问题：
\begin{itemize}
  \item 未来世界会怎样变化？
  \item 哪些动作与这些未来变化相匹配？
\end{itemize}

这类方法强调动作和视频之间的同步建模。动作不是附加模块，而是和未来视觉 token 一起进入模型。

\subsection{训练目标}

训练数据来自机器人演示轨迹。每个训练样本通常包含过去观测、语言指令、未来观测和未来动作：
\[
  (o_{\le t},\ell,o_{t+1:t+H},a_{t:t+H-1}).
\]
Joint Video-Action Modeling 的目标是让模型在给定过去观测和语言指令时，同时恢复未来视觉内容和动作序列。

如果使用扩散模型，可以令干净数据为：
\[
  x_0=[z^o_{t+1:t+H},z^a_{t:t+H-1}],
\]
其中 $z^o$ 是未来视频或图像的 latent token，$z^a$ 是动作 token。这里的 $x_0$ 不是单纯图像，而是“未来世界 token + 动作 token”的拼接对象。条件为：
\[
  c=[o_{\le t},\ell].
\]
训练时先随机采样一个噪声强度 $\tau$，再把高斯噪声 $\epsilon$ 加到干净样本 $x_0$ 上，得到带噪样本：
\[
  x_\tau=\sqrt{\bar{\alpha}_\tau}x_0+
  \sqrt{1-\bar{\alpha}_\tau}\epsilon,
  \qquad
  \epsilon\sim\mathcal{N}(0,I).
\]
模型看到的是 $(x_\tau,\tau,c)$，任务是预测刚才加入的噪声 $\epsilon$。扩散训练目标为：
\[
  \mathcal{L}_{\mathrm{joint}}
  =
  \mathbb{E}_{x_0,\epsilon,\tau}
  \left[
    \left\|
      \epsilon-\epsilon_\theta(x_\tau,\tau,c)
    \right\|_2^2
  \right].
\]
这个 loss 的含义是：模型必须学会在条件 $c$ 的约束下，把一个被污染的“未来世界+动作”样本逐步去噪回真实演示数据。训练完成后，推理时可以从纯噪声开始，反复去噪，生成与当前观测和语言指令一致的未来视觉 token 与动作 token。

如果未来视觉、动作和 proprioception 的尺度差异很大，实践中常会对不同 token 使用不同编码器、归一化方式或 loss 权重：
\[
  \mathcal{L}_{\mathrm{joint}}
  =
  \lambda_o\mathcal{L}_{\mathrm{video}}
  +
  \lambda_a\mathcal{L}_{\mathrm{action}}
  +
  \lambda_q\mathcal{L}_{\mathrm{state}}.
\]
但概念上，它们仍然服务于同一个目标：学习一个条件联合分布，使动作和未来世界在同一个生成过程中互相约束。

\subsection{推理流程}

推理时，模型从噪声开始生成：
\[
  (\hat{o}_{t+1:t+H},\hat{a}_{t:t+H-1})
  \sim
  p_\theta(\cdot\mid o_{\le t},\ell).
\]
执行时通常只执行第一个动作或 action chunk：
\[
  a_t=\hat{a}_t
  \quad \text{或} \quad
  a_{t:t+K-1}=\hat{a}_{t:t+K-1}.
\]

\subsection{优点与问题}

优点是结构统一，动作和未来观测可以互相约束。模型生成的动作天然对应某种未来世界变化，因此适合表达多模态动作分布和复杂交互。

问题是推理成本较高，尤其当未来视频 token 很多时，扩散去噪会非常慢。此外，联合建模要求模型同时处理视觉、动作和语言，训练难度较高。

\section{范式二：Imagine-Then-Execute}

\subsection{核心思想}

Imagine-Then-Execute 把 WAM 分成两个阶段：
\[
  \text{先想象未来，再根据未来选择动作。}
\]
第一阶段生成未来观测：
\[
  \hat{o}_{t+1:t+H}
  \sim
  p_\theta(o_{t+1:t+H}\mid o_{\le t},\ell).
\]
第二阶段根据当前观测和想象出的未来预测动作：
\[
  a_{t:t+H-1}
  \sim
  \pi_\phi(a_{t:t+H-1}\mid o_{\le t},\hat{o}_{t+1:t+H},\ell).
\]

它的直觉是：如果模型能想象任务成功的未来状态，那么动作模型可以学习如何到达这个未来。

\subsection{与 inverse dynamics 的关系}

第二阶段常常接近 inverse dynamics problem。Forward dynamics 问：
\[
  (o_t,a_t)\mapsto o_{t+1}.
\]
Inverse dynamics 问：
\[
  (o_t,o_{t+1})\mapsto a_t.
\]
Imagine-Then-Execute 可以看作先生成目标未来，再用 inverse dynamics 或 action decoder 反推出动作。

\subsection{训练目标}

Imagine-Then-Execute 通常把训练拆成两个目标。第一部分训练 future model，使其根据过去观测和语言指令生成未来观测：
\[
  \mathcal{L}_{\mathrm{future}}
  =
  -\log p_\theta(o_{t+1:t+H}\mid o_{\le t},\ell).
\]
如果 future model 是扩散视频模型，这个负对数似然通常也通过 denoising loss 近似：模型学习从带噪未来视频 token 中预测噪声。

第二部分训练 action model，使其根据当前观测和未来观测恢复动作：
\[
  \mathcal{L}_{\mathrm{action}}
  =
  -\log \pi_\phi(a_{t:t+H-1}\mid o_{\le t},o_{t+1:t+H},\ell).
\]
在训练时，这里的未来观测通常使用真实演示中的未来帧，而不是模型生成的未来帧。这样 action model 学到的是：如果当前状态和目标未来都给定，哪些动作能连接二者。

总训练目标可以写成：
\[
  \mathcal{L}
  =
  \mathcal{L}_{\mathrm{future}}
  +
  \beta\mathcal{L}_{\mathrm{action}}.
\]
因此，这一范式的训练逻辑是“先学会想象未来，再学会从未来反推动作”。它的风险在于训练时 action model 看到的未来是真实未来，推理时看到的未来却是生成未来，两者分布可能不完全一致。

\subsection{Planning 形式}

如果生成多个未来候选：
\[
  \hat{o}^{(i)}_{t+1:t+H}\sim p_\theta,
  \qquad i=1,\ldots,N,
\]
可以用 reward model 或 value model 评分：
\[
  J^{(i)}=\hat{V}_\psi(\hat{o}^{(i)}_{t+H}).
\]
选择：
\[
  i^*=\arg\max_i J^{(i)}.
\]
再由动作模型生成对应动作：
\[
  a_{t:t+H-1}
  \sim
  \pi_\phi(\cdot\mid o_{\le t},\hat{o}^{(i^*)}_{t+1:t+H},\ell).
\]

\subsection{优点与问题}

优点是模块化清晰：视频模型负责未来想象，动作模型负责执行。这种结构便于利用大规模预训练视频模型，也便于在未来空间中做候选选择。

问题是显式 future imagination 很慢，尤其是视频扩散模型需要多步 denoising。另一个问题是 generated future 和 executable action 之间可能不一致：模型可能想象出看起来合理但机器人实际上难以实现的未来。

\section{范式三：Fast-WAM}

\subsection{核心思想}

Fast-WAM 的出发点是质疑一个常见假设：WAM 是否必须在测试时显式生成未来视频？

它将训练和推理解耦。训练时保留视频预测目标：
\[
  \mathcal{L}
  =
  \mathcal{L}_{\mathrm{action}}
  +
  \lambda\mathcal{L}_{\mathrm{video}}.
\]
但推理时不生成未来视频，而是直接根据当前观测的 latent representation 输出动作：
\[
  a_{t:t+H-1}
  \sim
  \pi_\theta(a_{t:t+H-1}\mid z_t,\ell),
  \qquad
  z_t=f_\theta(o_{\le t}).
\]

\subsection{训练阶段}

训练时模型仍然学习 future prediction：
\[
  p_\theta(o_{t+1:t+H}\mid o_{\le t},a_{t:t+H-1},\ell),
\]
或在 latent space 中预测未来视频 token。该目标迫使模型学习物理动态、运动规律和任务相关时间结构。

同时，模型学习动作预测：
\[
  p_\theta(a_{t:t+H-1}\mid o_{\le t},\ell).
\]
因此视频预测在这里主要是 representation learning 的辅助目标。

更具体地说，Fast-WAM 的训练目标可以拆成两部分。第一部分是动作预测：
\[
  \mathcal{L}_{\mathrm{action}}
  =
  -\log p_\theta(a_{t:t+H-1}\mid z_t,\ell),
  \qquad
  z_t=f_\theta(o_{\le t}).
\]
第二部分是视频预测或视频 denoising：
\[
  \mathcal{L}_{\mathrm{video}}
  =
  -\log p_\theta(o_{t+1:t+H}\mid z_t,a_{t:t+H-1},\ell),
\]
或其扩散形式。总目标为：
\[
  \mathcal{L}
  =
  \mathcal{L}_{\mathrm{action}}
  +
  \lambda\mathcal{L}_{\mathrm{video}}.
\]

这里最重要的是 $\mathcal{L}_{\mathrm{video}}$ 的角色变化。它不是为了让模型在测试时真的生成未来视频，而是作为辅助监督，迫使 $z_t$ 包含对物理动态、物体交互和任务进展有用的信息。换言之，Fast-WAM 把视频预测从“推理步骤”降级为“训练时的表征学习约束”。

\subsection{推理阶段}

推理时，模型不再执行：
\[
  \hat{o}_{t+1:t+H}\sim p_\theta.
\]
它只使用视频模型或 world encoder 提供的 latent state：
\[
  z_t=f_\theta(o_{\le t}),
\]
然后由 action head 或 action expert 生成动作：
\[
  \hat{a}_{t:t+H-1}=g_\theta(z_t,\ell).
\]

\subsection{关键假设}

Fast-WAM 的核心假设是：
\[
  \text{WAM 的主要收益来自训练时的视频预测目标，而不是测试时显式生成未来。}
\]
换句话说，未来预测的作用是塑造更好的世界表征，而不一定要在推理时真的把未来视频生成出来。

\subsection{优点与问题}

优点是推理速度快，避免了测试时视频扩散的高延迟。它更适合 real-time control，尤其是机器人任务中控制频率有限、延迟敏感。

问题是它牺牲了显式 planning 的能力。模型不能在测试时比较多个想象未来，也较难利用 value ranking 做 best-of-$N$ 搜索。因此它更像一个被 world-model training 强化过的 policy，而不是完整的 model-based planner。

\section{三种范式对比}

\begin{center}
\begin{tabular}{llll}
\toprule
范式 & 训练时是否预测未来 & 推理时是否生成未来 & 动作如何产生 \\
\midrule
Joint Video-Action & 是 & 是 & 与未来视频联合生成 \\
Imagine-Then-Execute & 是 & 是 & 由想象未来反推出动作 \\
Fast-WAM & 是 & 否 & 由 latent representation 直接预测 \\
\bottomrule
\end{tabular}
\end{center}

从建模角度看：
\[
  \text{Joint Video-Action}
  :
  p(o_{future},a_{future}\mid o_{past},\ell).
\]
\[
  \text{Imagine-Then-Execute}
  :
  p(o_{future}\mid o_{past},\ell)
  \quad \text{then} \quad
  p(a_{future}\mid o_{past},o_{future},\ell).
\]
\[
  \text{Fast-WAM}
  :
  \mathcal{L}_{video}\text{ for training},
  \quad
  p(a_{future}\mid z_t,\ell)\text{ for inference}.
\]

\section{如何阅读相关论文}

阅读一篇 WAM 论文时，应先判断它属于哪种范式。可以依次追问：
\begin{enumerate}
  \item 动作和未来视频是联合生成，还是分阶段生成？
  \item 视频预测只用于训练，还是也用于测试时决策？
  \item 推理时是否显式生成未来图像或视频？
  \item 是否使用 value function 对未来候选排序？
  \item 动作是单步输出，还是 action chunk？
  \item 延迟瓶颈来自视频生成、动作扩散，还是多候选采样？
\end{enumerate}

\section{与三篇材料的对应}

\begin{itemize}
  \item Fast-WAM 对应第三种范式。它的贡献在于证明或至少经验性支持：训练时的视频建模可能比测试时显式未来想象更关键。
  \item Cosmos Policy 更接近联合建模与 planning 的组合。它把 action、future state 和 value 放入统一生成框架，并用候选未来与价值预测辅助决策。
  \item WM as simulator 不是典型 WAM 入门论文。它更关注把 world model 作为 learned simulator，用来生成训练场景或 synthetic data。
\end{itemize}

\section{一句话总结}

Joint Video-Action 的核心是“动作和未来一起生成”；Imagine-Then-Execute 的核心是“先想象未来，再反推动作”；Fast-WAM 的核心是“训练时学习未来，推理时直接行动”。

三者的根本分歧在于：未来预测到底应该作为测试时 planning 的显式步骤，还是只作为训练时塑造世界表征的辅助目标。

\end{document}
